% Guia do professor — Capítulo 23 (módulo extra): Mudança Climática
\section{Capítulo 23 (módulo extra) --- Mudança Climática}

\subsection*{Visão geral e ritmo}

O fecho do curso: nada de novo entra — estufa, feedbacks, oceano e
estatística já foram construídos; aqui são aplicados à perturbação
antrópica com números na mesa. Tom recomendado: física serena e
quantitativa; deixar os dados falarem. \textbf{Ritmo: 2 aulas de 2\,h
(a última pode fechar com retrospectiva do curso).}

\begin{itemize}
  \item \textbf{Aula 1} — o logaritmo do CO$_2$ (T1 no quadro),
        sensibilidade e nuvens, a inércia das duas caixas (Casos~1--2).
  \item \textbf{Aula 2} — TCRE e orçamento (Caso~3), detecção no
        ruído (Caso~4), impressões digitais (estratosfera!), impactos
        capítulo a capítulo e o fecho do curso.
\end{itemize}

\subsection*{Objetivos de aprendizagem}

\begin{enumerate}
  \item deduzir/justificar $\Delta F = 5{,}35\ln(C/C_0)$ e calcular a
        forçante atual;
  \item combinar forçante e feedbacks para a faixa de sensibilidade;
  \item explicar o aquecimento comprometido pelas escalas oceânicas;
  \item usar o TCRE para converter metas em orçamentos de carbono;
  \item distinguir sinal de ruído e listar as assinaturas do
        mecanismo estufa.
\end{enumerate}

\subsection*{Pré-requisitos a reativar}

Bandas e janelas (Cap.~2); feedbacks e $\lambda_0$ (Cap.~21);
capacidade térmica do oceano (Cap.~3); ENSO como ruído (Cap.~21);
tendência vs.\ variância (estatística básica); MPI $+2$\,K (Cap.~16,
N3).

\subsection*{Armadilhas conceituais frequentes}

\begin{itemize}
  \item \textbf{``A banda está saturada, mais CO$_2$ não faz nada''}:
        o argumento do centro ignora as asas E a altura de emissão —
        T1 e N1 desmontam com números.
  \item \textbf{``1 K não é nada''}: 1\,K GLOBAL move distribuições
        inteiras — os extremos (Cap.~4: $+7\%$/K) escalam muito mais
        rápido que a média.
  \item \textbf{``Uma década fria refuta''}: detecção é multidecadal
        (Caso~4); ruído interno esconde o sinal por 30--40 anos.
  \item \textbf{``É o Sol''}: a estratosfera ESFRIA (T5) — o Sol
        aqueceria tudo junto; a assinatura vertical decide.
  \item \textbf{``Estabilizar emissões estabiliza o clima''}: TCRE
        exige líquido ZERO; estabilizar emissões só estabiliza a
        TAXA de aquecimento.
\end{itemize}

\subsection*{Demonstração de baixo custo}

Espectro de OLR medido por satélite (IRIS/AIRS) com o ``dente'' do
CO$_2$ — a estufa fotografada. Série do Niño-3.4 sobre a série global
de temperatura: o ruído do Caso~4 na vida real. Keeling completo
(dente de serra sazonal $+$ tendência): dois fenômenos do curso num
gráfico.

\subsection*{Problemas sugeridos do livro (módulo extra — sem seção correspondente no livro; usar os complementares)}

Aquecimento: $\Delta F$ para concentrações dadas; aquecimento de
equilíbrio por sensibilidade. Aprofundamento: orçamentos para metas;
expansão térmica. Síntese: ``liste três previsões do mecanismo estufa
feitas ANTES de serem observadas'' (estratosfera fria, noites
aquecendo mais, Ártico amplificado).
\emph{Confira a numeração na sua cópia (v1.02b).}

\subsection*{Uso do notebook \texttt{cap23\_mudanca.ipynb}}

\begin{center}
\begin{tabular}{lll}
\hline
Caso & Conteúdo & Momento sugerido \\
\hline
1 & forçante log e cenários por sensibilidade & Aula 1 \\
2 & duas caixas: aquecimento comprometido & Aula 1 \\
3 & TCRE e orçamento de carbono & Aula 2 \\
4 & emergência do sinal no ruído & Aula 2 \\
\hline
\end{tabular}
\end{center}

O Caso~2 congelando o CO$_2$ em 2024 rende a melhor discussão do
semestre: a dívida térmica como consequência direta do Cap.~3.

\subsection*{Exercícios complementares e soluções}

Nas notas: T1--T5 e N1--N4. Soluções em \texttt{cap23\_solucoes.ipynb}
(\emph{não distribuir}): T2 com \texttt{sympy} (escalas de tempo); N1
recupera o 5,35 de um modelo de banda; N3 compara trajetórias de
mesmo acumulado. Pares: T1$\leftrightarrow$N1, T2$\leftrightarrow$N2,
T4$\leftrightarrow$N3.

\subsection*{Conexões com o resto do curso}

TODAS, de propósito: estufa (Cap.~2), vapor e extremos (Cap.~4),
nuvens (Caps.~6--7), satélites que medem o dente do CO$_2$ (Cap.~8),
jato e bloqueios (Cap.~11), ciclones num mundo quente (Caps.~13, 16),
oceano lento (Cap.~3), modelos e ensembles (Cap.~20), clima natural
como régua (Cap.~21). O capítulo é a prova de que o curso forma
leitores qualificados do IPCC.
